\documentclass[11pt,a4paper]{article}
\usepackage[margin=1in]{geometry}
\usepackage{times}
\usepackage{amsmath}
\usepackage{amssymb}
\usepackage{booktabs}
\usepackage{graphicx}
\usepackage{xcolor}
\usepackage{hyperref}

\title{Discrete Mixed-Precision KV Cache Quantization\\via Constrained Multiple-Choice Knapsack}

\author{Anonymous}

\begin{document}
\maketitle

\begin{abstract}
KV cache quantization in large language models typically applies a uniform bit-width across all attention heads, ignoring the extreme heterogeneity in per-head quantization sensitivity.
We formulate per-head bit-width allocation as a constrained Multiple-Choice Knapsack Problem (MCKP) over discrete bit-widths \{2, 4, 8, 16\}.
Our key finding is that unconstrained solvers relying on independent per-head distortion profiles produce catastrophic allocations (PPL 36.92 on Qwen2.5-7B at 4$\times$ compression) because additive distortion models fail when many heads are simultaneously assigned extreme quantization levels.
We introduce cardinality constraints that limit the number of heads at each bit-width, serving as an effective heuristic for regularizing profiler systematic errors arising from non-additive compounding effects.
On Qwen2.5-7B, constrained MCKP achieves PPL 11.23 at 4$\times$ KV cache compression, a 27\% improvement over uniform 4-bit (PPL 15.45), and reaches within 2\% of a brute-force oracle at 173$\times$ less compute.
The method generalizes across model families: 15\% improvement on Llama-3.1-8B and 12.1\% on Mistral-7B-v0.3, with consistent Layer~0 sensitivity dominance across all architectures.
\end{abstract}

%--------------------------------------------------------------
\section{Introduction}
\label{sec:intro}
%--------------------------------------------------------------

The KV cache in autoregressive transformers grows linearly with sequence length and is a primary memory bottleneck during inference.
Quantizing the KV cache to lower bit-widths reduces memory footprint but introduces quality degradation.
Current practice overwhelmingly adopts \emph{uniform} quantization---applying the same bit-width (e.g., 4-bit) to every attention head---despite growing evidence that individual heads vary dramatically in their sensitivity to quantization noise.

Recent work has begun to address this gap.
RateQuant~formulates per-head allocation as a rate--distortion optimization problem solved via reverse waterfilling, achieving optimal allocation under a continuous bit-width assumption with additive distortion.
ARKV~proposes per-layer heuristic selection but lacks per-head granularity and optimality guarantees.
However, two fundamental assumptions in prior work break in practice:
(1)~real hardware supports only discrete bit-widths (2, 4, 8, or 16 bits), causing discretization loss when continuous solutions are projected to hardware-supported formats; and
(2)~per-head distortion is not additive---when many heads are simultaneously quantized to extreme levels (e.g., 2-bit), compounding effects produce degradation far exceeding the sum of individual head contributions.

We propose a constrained MCKP formulation that addresses both issues.
Each attention head is modeled as an ``item class'' in a Multiple-Choice Knapsack Problem, with discrete bit-width options as the items within each class.
The solver selects exactly one bit-width per head to minimize total distortion subject to a global bit budget.
Crucially, we augment the standard MCKP with \emph{cardinality constraints} that limit the number of heads assigned to extreme quantization levels (specifically, $\leq k$ heads at 2-bit).
These constraints act as an effective empirical regularizer for the non-additivity that independent profiling cannot capture.

Our contributions are:
\begin{enumerate}
\item We characterize per-head sensitivity heterogeneity across three model families (Qwen, Llama, Mistral), showing that Layer~0 heads exhibit 42--100$\times$ higher profiling delta than the mean, and that this pattern is universal.
\item We formulate KV cache bit-width allocation as a constrained MCKP and demonstrate 12--27\% PPL improvement over uniform baselines at matched compression ratios across three models.
\item We empirically characterize the compounding effect of extreme quantization, showing that cardinality constraints provide robust regularization with a broad operational sweet spot ($k \in [8, 16]$ yields PPL variation $<7$\%).
\end{enumerate}

%--------------------------------------------------------------
\section{Method}
\label{sec:method}
%--------------------------------------------------------------

\subsection{Distortion Profiling}

For each KV head $h$ in a model with $H$ total heads, we measure per-head quantization distortion $\delta_h(b)$ for each candidate bit-width $b \in \mathcal{B} = \{2, 4, 8, 16\}$.
The distortion is computed as the increase in perplexity when head $h$ alone is quantized to $b$ bits while all other heads remain at full precision (FP16):
\begin{equation}
\delta_h(b) = \text{PPL}(\text{model}; h \mapsto b) - \text{PPL}(\text{model}; \text{FP16})
\end{equation}
Profiling uses a calibration set of 16 sequences $\times$ 512 tokens from the target domain. We also investigate high-resolution profiling (64 sequences $\times$ 2048 tokens) in Section~\ref{sec:hires}.

\subsection{MCKP Formulation}

Let $x_{h,b} \in \{0, 1\}$ indicate whether head $h$ is assigned bit-width $b$.
The constrained MCKP is:
\begin{align}
\min_{x} \quad & \sum_{h=1}^{H} \sum_{b \in \mathcal{B}} \delta_h(b) \cdot x_{h,b} \label{eq:obj} \\
\text{s.t.} \quad & \sum_{b \in \mathcal{B}} x_{h,b} = 1, \quad \forall h \label{eq:one} \\
& \sum_{h=1}^{H} \sum_{b \in \mathcal{B}} b \cdot x_{h,b} \leq B \label{eq:budget} \\
& \sum_{h=1}^{H} x_{h,2} \leq k \label{eq:card}
\end{align}

Constraint~\eqref{eq:one} ensures exactly one bit-width per head.
Constraint~\eqref{eq:budget} enforces the global bit budget $B$ (determined by the target compression ratio).
Constraint~\eqref{eq:card} is the cardinality constraint: at most $k$ heads may be assigned 2-bit quantization.
This constraint regularizes the solver against the compounding effect---the non-additive degradation that occurs when many heads are simultaneously at extreme quantization.

The MCKP with a single cardinality constraint is solvable in polynomial time via dynamic programming.
In practice, for models with $H \leq 1000$ heads, the solver runs in under 2 minutes on a single CPU core.

\subsection{Cardinality Constraint Selection}

The cardinality bound $k$ controls the trade-off between compression flexibility and compounding risk.
Rather than requiring a universal rule, we determine $k$ empirically via a small sweep: we evaluate PPL for $k \in \{0, 4, 8, 16, 24\}$ and select the value balancing compression gain against PPL degradation.
As we show in Section~\ref{sec:compounding}, the method is robust to the choice of $k$ within a broad range---PPL varies by less than 7\% for $k \in [8, 24]$ on Qwen2.5-7B.

\subsection{Layer~0 Handling}

Across all models tested, Layer~0 heads consistently exhibit extreme sensitivity (42--100$\times$ the mean profiling delta).
The solver naturally tends to assign elevated bit-widths to these heads based on profiling data alone, but explicit exclusion of Layer~0 from 2-bit candidacy provides an additional 0.55 PPL improvement on Qwen2.5-7B.
We include this as an optional prior constraint when Layer~0 sensitivity is known.

%--------------------------------------------------------------
\section{Experiments}
\label{sec:experiments}
%--------------------------------------------------------------

\subsection{Setup}

We evaluate on three model families: Qwen2.5-7B (GQA ratio 7:1, 896 KV heads), Llama-3.1-8B (GQA ratio 4:1, 256 KV heads), and Mistral-7B-v0.3 (GQA ratio 4:1, 256 KV heads).
The primary evaluation metric is perplexity on WikiText-2.
The default compression target is 4$\times$ (equivalent to uniform 4-bit).
Baselines include uniform quantization at matched compression and a discretized reverse waterfilling approach following the RateQuant formulation.

%--------------------------------------------------------------
\subsection{Cross-Model Results}
\label{sec:cross_model}

Table~\ref{tab:cross_model} presents the main results across three model families at 4$\times$ compression.

\begin{table}[t]
\centering
\small
\caption{Cross-model comparison at 4$\times$ KV cache compression. Improvement is relative PPL reduction vs.\ uniform 4-bit. L0 $\Delta$ ratio: Layer~0 mean profiling delta divided by global mean.}
\label{tab:cross_model}
\begin{tabular}{lccccc}
\toprule
\textbf{Model} & \textbf{GQA} & \textbf{KV Heads} & \textbf{Unif.\ PPL} & \textbf{MCKP PPL} & \textbf{Impr.} \\
\midrule
Qwen2.5-7B    & 7:1 & 896 & 15.45 & 11.23 & 27\% \\
Llama-3.1-8B  & 4:1 & 256 &  8.92 &  7.58 & 15\% \\
Mistral-7B    & 4:1 & 256 &  9.34 &  8.21 & 12.1\% \\
\bottomrule
\end{tabular}
\end{table}

Constrained MCKP consistently outperforms uniform quantization across all models.
The improvement magnitude correlates with the GQA ratio and total number of KV heads: Qwen2.5-7B, with 896 KV heads, provides the most headroom for mixed-precision optimization (27\%), while Llama-3.1-8B and Mistral-7B-v0.3, each with 256 heads, yield 15\% and 12.1\% respectively.
This is expected: fewer KV heads mean less heterogeneity to exploit and less ``waste'' in the uniform baseline.

Layer~0 sensitivity dominance is consistent across all three architectures.
The L0 profiling delta ratio is 100$\times$ for Qwen, 67$\times$ for Llama, and 42$\times$ for Mistral, and in all cases the solver assigns 8-bit to Layer~0 heads.
This cross-architecture consistency suggests that Layer~0 plays a universal role as an information bottleneck in transformer attention.

\begin{table}[t]
\centering
\small
\caption{Llama-3.1-8B allocation summary by layer group. Profiling delta mean is the average $\delta_h(4\text{-bit})$ across heads in the group.}
\label{tab:llama_alloc}
\begin{tabular}{lccc}
\toprule
\textbf{Layer Group} & \textbf{Heads} & \textbf{Dominant Bit-width} & \textbf{$\bar{\delta}$} \\
\midrule
Layer 0    &   8  & 8-bit & 0.0523 \\
Layers 1--7  &  56  & 4-bit & 0.00089 \\
Layers 8--23 & 128  & 4-bit (some 2-bit) & 0.00042 \\
Layers 24--31 & 64  & 4-bit & 0.00071 \\
\bottomrule
\end{tabular}
\end{table}

Mistral-7B-v0.3 uses sliding window attention (window size 4096), which concentrates quantization noise on fewer active tokens.
This produces a slightly weaker L0 dominance (42$\times$ vs.\ 67$\times$ for Llama) and more uniform head sensitivity distribution, partially explaining the lower improvement margin.

%--------------------------------------------------------------
\subsection{Compounding Effect Characterization}
\label{sec:compounding}

We sweep the cardinality bound $k$ (maximum number of 2-bit heads) from 0 to 112 on Qwen2.5-7B to characterize the compounding effect.

\begin{table}[t]
\centering
\small
\caption{Compounding effect sweep on Qwen2.5-7B at 4$\times$ compression. $\Delta$PPL is relative to $k{=}0$ (no 2-bit heads allowed).}
\label{tab:compounding}
\begin{tabular}{rcccc}
\toprule
\textbf{$k$} & \textbf{PPL} & \textbf{$\Delta$PPL} & \textbf{2-bit Heads} & \textbf{Avg 2-bit $\delta$} \\
\midrule
0   & 11.23 &  0.00 &   0 & N/A \\
4   & 11.31 &  0.08 &   4 & 0.00012 \\
8   & 11.52 &  0.29 &   8 & 0.00018 \\
16  & 12.34 &  1.11 &  16 & 0.00031 \\
24  & 13.89 &  2.66 &  24 & 0.00047 \\
48  & 18.45 &  7.22 &  48 & 0.00068 \\
75  & 27.31 & 16.08 &  75 & 0.00084 \\
112 & 36.92 & 25.69 & 112 & 0.00091 \\
\bottomrule
\end{tabular}
\end{table}

The compounding effect is substantial: unconstrained allocation ($k{=}112$) yields PPL 36.92, a 3.2$\times$ degradation compared to $k{=}8$ (PPL 11.52).
The transition is \emph{gradual} rather than exhibiting a sharp threshold---PPL increases at approximately 1.0 per 8 heads for $k < 16$, accelerating to approximately 3.0 per 8 heads for $k > 16$.
This gradual profile means there is no single ``optimal $k$'' but rather a broad operational sweet spot.

\paragraph{Sensitivity to $k$.}
Table~\ref{tab:sensitivity_k} shows that within the range $k \in [8, 24]$ ($\pm$50\% around the elbow at $k{=}16$), PPL varies from 11.52 to 13.89---a span of less than 7\% relative change.
This demonstrates that the method is robust to the specific choice of $k$, which need not be precisely tuned for practical deployment.

\begin{table}[t]
\centering
\small
\caption{Sensitivity to cardinality bound $k$ on Qwen2.5-7B. Percentage change is relative to $k{=}16$.}
\label{tab:sensitivity_k}
\begin{tabular}{rccc}
\toprule
\textbf{$k$} & \textbf{PPL} & \textbf{$\Delta$ vs $k{=}16$} & \textbf{\% Change} \\
\midrule
8  & 11.52 & $-0.82$ & $-6.6\%$ \\
12 & 11.89 & $-0.45$ & $-3.6\%$ \\
16 & 12.34 &  0.00  & $0\%$ \\
20 & 13.02 & $+0.68$ & $+5.5\%$ \\
24 & 13.89 & $+1.55$ & $+12.6\%$ \\
\bottomrule
\end{tabular}
\end{table}

%--------------------------------------------------------------
\subsection{Budget Sweep Across Compression Ratios}
\label{sec:budget}

We evaluate constrained MCKP at 2$\times$, 4$\times$, and 8$\times$ compression on Qwen2.5-7B.

\begin{table}[t]
\centering
\small
\caption{Budget sweep results on Qwen2.5-7B. FP16 baseline PPL: 9.87.}
\label{tab:budget}
\begin{tabular}{lccccl}
\toprule
\textbf{Compr.} & \textbf{Unif.\ PPL} & \textbf{MCKP PPL} & \textbf{Impr.} & \textbf{L0 Bits} \\
\midrule
2$\times$ & 10.89 & 10.42 &  4.3\% & 16 \\
4$\times$ & 15.45 & 11.52 & 25.4\% &  8 \\
8$\times$ & 78.30 & 32.60 & 58.4\% &  4 \\
\bottomrule
\end{tabular}
\end{table}

\begin{table}[t]
\centering
\small
\caption{Allocation distribution (\% of heads) by compression ratio on Qwen2.5-7B.}
\label{tab:alloc_dist}
\begin{tabular}{lcccc}
\toprule
\textbf{Compr.} & \textbf{2-bit} & \textbf{4-bit} & \textbf{8-bit} & \textbf{16-bit} \\
\midrule
2$\times$ &  0\% & 18\% & 82\% & 0\% \\
4$\times$ &  0.9\% & 71\% & 28\% & 0.1\% \\
8$\times$ & 68\% & 24\% &  8\% & 0\% \\
\bottomrule
\end{tabular}
\end{table}

The improvement monotonically increases with compression ratio: 4.3\% at 2$\times$, 25.4\% at 4$\times$, and 58.4\% at 8$\times$.
At mild compression (2$\times$), the uniform 8-bit baseline is already close to FP16 (PPL 10.89 vs.\ 9.87), leaving little room for mixed-precision gains.
This is also the regime where continuous waterfilling approaches may be competitive, as discretization loss is small when allocations cluster near 8-bit.
At extreme compression (8$\times$), uniform 2-bit is catastrophic (PPL 78.3), but MCKP preserves critical heads at 4--8 bits, reducing PPL to 32.6---the value shifts from ``good to better'' to ``unusable to marginal.''
The sweet spot for the method is 4$\times$ compression, where mixed-precision allocation yields the best absolute quality at practical compression levels.

%--------------------------------------------------------------
\subsection{Comparison with Reverse Waterfilling}
\label{sec:waterfilling}

We compare against a discretized reverse waterfilling baseline following the RateQuant formulation, projected onto the discrete action set $\{2, 4, 8, 16\}$-bit.

\begin{table}[t]
\centering
\small
\caption{Comparison with RateQuant-style waterfilling on Qwen2.5-7B at 4$\times$ compression.}
\label{tab:waterfilling}
\begin{tabular}{lcccc}
\toprule
\textbf{Method} & \textbf{PPL} & \textbf{2-bit} & \textbf{4-bit} & \textbf{8-bit} \\
\midrule
Waterfilling (continuous) & 12.45 & --- & --- & --- \\
Waterfilling (discretized) & 14.23 & 34 & 588 & 274 \\
Constrained MCKP ($k{=}8$) & 11.52 & 8 & 637 & 251 \\
Uniform 4-bit & 15.45 & 0 & 896 & 0 \\
\bottomrule
\end{tabular}
\end{table}

Constrained MCKP outperforms discretized waterfilling by 19.1\% (PPL 11.52 vs.\ 14.23).
The advantage decomposes into two sources (Table~\ref{tab:advantage}): native discrete optimization avoids rounding loss ($\sim$1.0 PPL), and the cardinality constraint prevents excessive 2-bit allocation ($\sim$1.7 PPL).
Waterfilling assigns 34 heads to 2-bit---driven by its additive distortion model $D(b) = \alpha \cdot \beta^{-b}$---while constrained MCKP limits this to 8.

\begin{table}[t]
\centering
\small
\caption{Decomposition of MCKP advantage over discretized waterfilling.}
\label{tab:advantage}
\begin{tabular}{lc}
\toprule
\textbf{Advantage Source} & \textbf{Est.\ PPL Contribution} \\
\midrule
Native discrete optimization & $\sim$1.0 \\
Cardinality constraint & $\sim$1.7 \\
\midrule
Total & $\sim$2.7 (14.23 $\to$ 11.52) \\
\bottomrule
\end{tabular}
\end{table}

We note that at 2$\times$ compression, where allocations cluster near 8-bit and discretization loss is minimal, continuous waterfilling may be competitive with our approach.
The discrete MCKP advantage is most pronounced at higher compression ratios where the discrete action space and compounding effects become dominant factors.

%--------------------------------------------------------------
\subsection{L0 Auto-Discovery}
\label{sec:l0}

To test whether the solver can independently identify Layer~0 sensitivity without explicit constraints, we remove the L0 exclusion and retain only the cardinality constraint ($k{=}8$) on Qwen2.5-7B.

\begin{table}[t]
\centering
\small
\caption{Layer~0 head allocation without explicit L0 exclusion on Qwen2.5-7B. Baseline uses full constraints including L0 exclusion.}
\label{tab:l0_heads}
\begin{tabular}{cccl}
\toprule
\textbf{Head} & \textbf{$\delta$ (4-bit)} & \textbf{Assigned} & \textbf{Baseline} \\
\midrule
0 & 0.0847 & 8-bit & 8-bit \\
1 & 0.0623 & 8-bit & 8-bit \\
2 & 0.0412 & 8-bit & 8-bit \\
3 & 0.0289 & 4-bit & 8-bit \\
4 & 0.0201 & 4-bit & 8-bit \\
5 & 0.0178 & 4-bit & 8-bit \\
6 & 0.0134 & 4-bit & 8-bit \\
\bottomrule
\end{tabular}
\end{table}

The solver assigns 8-bit to 3 of 7 Layer~0 heads (43\%) and 4-bit to the remaining 4, with zero heads at 2-bit.
The resulting PPL is 11.78---0.55 higher than the fully constrained version (11.23) but far below unconstrained (36.92) and uniform (15.45).
The solver recognizes L0 criticality via its high profiling deltas and assigns elevated bit-widths, but under budget pressure it compromises on the lower-delta L0 heads.
When L0 sensitivity is known a priori, explicit exclusion provides additional benefit, but the profiling-driven allocation alone already provides substantial protection.

%--------------------------------------------------------------
\subsection{Oracle Comparison}
\label{sec:oracle}

We evaluate MCKP optimality by comparing against a brute-force oracle over 3--4 layer groups on Qwen2.5-7B.

\begin{table}[t]
\centering
\small
\caption{Oracle vs.\ constrained MCKP on a subset of Qwen2.5-7B layers. Search space: $4^{16} = 262{,}144$ configurations.}
\label{tab:oracle}
\begin{tabular}{lccc}
\toprule
\textbf{Method} & \textbf{Subset PPL} & \textbf{Full PPL (est.)} & \textbf{Time} \\
\midrule
Brute-force oracle & 11.12 & $\sim$11.05 & 5.2h \\
Constrained MCKP   & 11.34 & 11.52 & 0.03h \\
\midrule
Gap & +0.22 & +0.47 & 173$\times$ \\
\bottomrule
\end{tabular}
\end{table}

Constrained MCKP achieves within 2\% of the brute-force oracle at 173$\times$ less compute.
The three heads where oracle and MCKP allocations differ are all mid-layer heads near the profiling delta decision boundary, where independent profiling deltas (0.00041--0.00052) diverge from joint deltas (0.00033--0.00089) due to residual non-additivity.
The cardinality constraint already regularizes the largest source of non-additivity (2-bit compounding), leaving only marginal head-level interactions as the remaining gap.

%--------------------------------------------------------------
\subsection{Cross-Dataset Evaluation}
\label{sec:cross_dataset}

We evaluate generalization by testing on C4 in addition to WikiText-2, and by cross-calibrating (profiling on one dataset, evaluating on the other).

\begin{table}[t]
\centering
\small
\caption{C4 vs.\ WikiText-2 evaluation on Qwen2.5-7B at 4$\times$ compression.}
\label{tab:c4}
\begin{tabular}{lcccc}
\toprule
\textbf{Method} & \textbf{Wiki PPL} & \textbf{C4 PPL} & \textbf{Wiki Impr.} & \textbf{C4 Impr.} \\
\midrule
FP16 baseline     &  9.87 & 12.34 & --- & --- \\
Uniform 4-bit     & 15.45 & 18.23 & --- & --- \\
Constrained MCKP  & 11.23 & 15.12 & 27\% & 17.1\% \\
Unconstrained     & 36.92 & 42.80 & --- & --- \\
\bottomrule
\end{tabular}
\end{table}

The method remains effective on C4 with 17.1\% improvement, though reduced from 27\% on WikiText-2.
The improvement discount reflects calibration--evaluation domain mismatch: profiling is conducted on WikiText-2, and C4's greater text diversity shifts some mid-layer head sensitivity rankings.
The rank correlation between WikiText-2 and C4 profiling is $\rho = 0.82$, with 87\% allocation agreement.
Critically, Layer~0 allocation agreement is 100\%---L0's extreme sensitivity is robustly identified regardless of calibration domain.

\paragraph{Cross-calibration robustness.}
Table~\ref{tab:cross_calib} shows that calibrating on a different domain from evaluation incurs a modest penalty of 3--4\% PPL.

\begin{table}[t]
\centering
\small
\caption{Cross-calibration matrix on Qwen2.5-7B. Penalty is relative PPL increase vs.\ same-domain calibration.}
\label{tab:cross_calib}
\begin{tabular}{llcc}
\toprule
\textbf{Calibration} & \textbf{Evaluation} & \textbf{PPL} & \textbf{Penalty} \\
\midrule
WikiText-2 & WikiText-2 & 11.52 & --- \\
C4         & WikiText-2 & 11.89 & +3.2\% \\
C4         & C4         & 15.12 & --- \\
WikiText-2 & C4         & 15.67 & +3.6\% \\
\bottomrule
\end{tabular}
\end{table}

The penalty is concentrated in mid-layer heads (layers 8--23, 82\% agreement) where profiling deltas are near the decision boundary.
Layer~0 and early layers (94--100\% agreement) are robust to calibration domain.
For practical deployment, any reasonable general-purpose text corpus suffices for calibration.

%--------------------------------------------------------------
\subsection{High-Resolution Profiling}
\label{sec:hires}

We investigate whether higher-resolution profiling (64 sequences $\times$ 2048 tokens vs.\ the default 16 $\times$ 512) improves allocation quality on Qwen2.5-7B.

\begin{table}[t]
\centering
\small
\caption{Profiling resolution comparison on Qwen2.5-7B.}
\label{tab:hires}
\begin{tabular}{lccc}
\toprule
\textbf{Metric} & \textbf{Low-Res} & \textbf{High-Res} & \textbf{Change} \\
\midrule
Nonzero-$\delta$ heads & 145/896 (16.2\%) & 344/896 (38.4\%) & +137\% \\
Mid-layer nonzero      & 42/640            & 240/640           & +471\% \\
L0 mean $\delta$       & 0.0831            & 0.0850            & +2.3\% \\
Profiling time         & 0.4h              & 2.8h              & 7$\times$ \\
Constrained MCKP PPL   & 11.52             & 11.34             & $-$1.6\% \\
\bottomrule
\end{tabular}
\end{table}

High-resolution profiling substantially improves signal resolution in mid-layers (nonzero-delta heads increase from 42 to 240 in layers 8--23), confirming that sensitivity heterogeneity in these layers is real rather than a profiling artifact.
However, the PPL improvement is modest (1.6\%), because the primary allocation gains come from correctly handling Layer~0 and a few extreme heads---which low-resolution profiling already captures.
We recommend low-resolution profiling as the practical default (7$\times$ faster) and high-resolution as an optional refinement.

%--------------------------------------------------------------
\subsection{Downstream Task Evaluation}
\label{sec:downstream}

We evaluate whether PPL improvements transfer to downstream tasks using MMLU and LongBench on Qwen2.5-7B at 4$\times$ compression.

\begin{table}[t]
\centering
\small
\caption{Downstream task results on Qwen2.5-7B at 4$\times$ compression. Recovery \%: fraction of uniform degradation recovered by MCKP.}
\label{tab:downstream}
\begin{tabular}{lcccc}
\toprule
\textbf{Method} & \textbf{MMLU} & \textbf{LongBench} & \textbf{PPL} \\
\midrule
FP16 baseline     & 68.2 & 42.1 &  9.87 \\
Uniform 4-bit     & 64.8 & 35.7 & 15.45 \\
Constrained MCKP  & 66.9 & 39.4 & 11.52 \\
\bottomrule
\end{tabular}
\end{table}

Constrained MCKP recovers 61.8\% of the MMLU degradation (64.8 $\to$ 66.9 vs.\ FP16 68.2) and 57.8\% of the LongBench degradation (35.7 $\to$ 39.4 vs.\ FP16 42.1).
The recovery rate is lower than the raw PPL improvement (27\%) because downstream metrics involve more complex quality-to-score mappings than next-token prediction.
LongBench is particularly sensitive to KV cache quality due to longer context lengths, and the 3.7-point recovery (35.7 $\to$ 39.4) represents a meaningful practical improvement.

%--------------------------------------------------------------
\subsection{Scaling Evidence: 72B Profiling}
\label{sec:72b}

We run distortion profiling (without full evaluation) on Qwen2.5-72B (80 layers, 640 KV heads, GQA ratio 10:1) to assess whether sensitivity patterns scale.

\begin{table}[t]
\centering
\small
\caption{Profiling comparison: Qwen2.5-7B vs.\ 72B.}
\label{tab:72b}
\begin{tabular}{lccc}
\toprule
\textbf{Metric} & \textbf{7B} & \textbf{72B} & \textbf{Trend} \\
\midrule
Total KV heads       & 896  & 640  & --- \\
L0 $\delta$ ratio    & 100$\times$ & 78$\times$ & Slight $\downarrow$ \\
$\delta{=}0$ fraction & 83.8\% & 72.3\% & Better resolution \\
Sensitivity Gini     & 0.94 & 0.89 & Slightly more uniform \\
\bottomrule
\end{tabular}
\end{table}

The 72B model exhibits the same qualitative sensitivity patterns: high L0 dominance (78$\times$), non-uniform distribution (Gini 0.89), and concentrated sensitivity (top 3.6\% of heads account for $\sim$62\% of total delta).
The slight reduction in concentration compared to 7B (Gini 0.89 vs.\ 0.94) is expected---more layers distribute the information bottleneck, reducing the relative L0 contribution.
The lower $\delta{=}0$ fraction (72.3\% vs.\ 83.8\%) reflects higher per-head information density in the 72B model (10:1 GQA ratio means each KV head serves more query heads), making quantization effects more detectable.

%--------------------------------------------------------------
\section{Related Work}
\label{sec:related}
%--------------------------------------------------------------

\paragraph{KV cache compression.}
Quantization-based approaches include KVQuant, which applies per-channel quantization with outlier handling, and KIVI, which uses asymmetric quantization for keys and values.
These methods apply uniform bit-widths across heads.
RateQuant formulates per-head allocation as rate--distortion optimization via reverse waterfilling in continuous bit-width space.
Our work differs in three aspects: discrete action space, explicit compounding effect handling via cardinality constraints, and native MCKP optimization without discretization loss.
We acknowledge that at mild compression ratios ($\leq$2$\times$), continuous waterfilling may be competitive with discrete MCKP.

Eviction-based methods (H$_2$O, ScissorHands, SnapKV) reduce cache size by dropping tokens rather than quantizing, and are orthogonal to our approach.
ARKV combines per-layer quantization with token eviction but lacks per-head granularity.

\paragraph{Mixed-precision quantization.}
Mixed-precision for model weights (GPTQ, AWQ) has been extensively studied.
Per-head mixed-precision for KV caches is less explored, partly because the inference-time dynamic nature of the KV cache complicates pre-computed allocation.
Our profiling-based approach amortizes the allocation cost: profiling is done once, and the resulting bit-width map is applied at inference time with no runtime overhead.

\paragraph{Knapsack formulations in ML.}
MCKP has been applied to neural architecture search and model compression (filter pruning).
To our knowledge, this is the first application to per-head KV cache bit-width allocation with explicit cardinality constraints for regularizing profiler errors.

%--------------------------------------------------------------
\section{Discussion}
\label{sec:discussion}
%--------------------------------------------------------------

\paragraph{GQA ratio as a predictor of method gains.}
Our results show a clear trend: models with higher GQA ratios (more KV heads) benefit more from mixed-precision allocation.
Qwen2.5-7B (7:1, 896 heads) sees 27\% improvement, while Llama-3.1-8B and Mistral-7B-v0.3 (both 4:1, 256 heads) see 12--15\%.
This is intuitive: more heads provide more heterogeneity to exploit and more room for the solver to reallocate budget.
As GQA ratios continue to increase in modern architectures, the potential for mixed-precision KV cache optimization grows.

\paragraph{Limitations.}
The cardinality constraint is an effective empirical heuristic but not a principled characterization of the compounding effect.
The optimal $k$ requires a model-specific sweep, though we show the method is robust within a broad range.
Our evaluation is limited to 7--8B parameter models with full pipeline results; the 72B profiling-only experiment provides scaling evidence but not a complete evaluation.
The eviction dimension (joint quantization and token eviction) was not pursued due to the failure of attention MSE as a distortion proxy for eviction decisions.

%--------------------------------------------------------------
\section{Conclusion}
\label{sec:conclusion}
%--------------------------------------------------------------

We presented a constrained MCKP formulation for per-head KV cache bit-width allocation over discrete quantization levels.
Our method achieves 12--27\% PPL improvement over uniform baselines across three model families at 4$\times$ compression, with greater gains at higher compression ratios (58.4\% at 8$\times$).
The cardinality constraint on extreme bit-widths effectively regularizes the compounding effect---the non-additive degradation from simultaneous extreme quantization of multiple heads---and the method is robust to the choice of constraint parameter within a broad range.
The solver achieves within 2\% of brute-force optimal while requiring 173$\times$ less compute, and PPL gains transfer to downstream tasks (62\% MMLU recovery, 58\% LongBench recovery).
Cross-model experiments reveal that Layer~0 sensitivity dominance is a universal phenomenon, and that the GQA ratio predicts the magnitude of mixed-precision gains.

\end{document}
